\section{Related Work}
A well-studied problem of AI in team formation is the \textit{Hedonic Games} framework~\cite{aziz2017fractional, gairing2010computing}, %, ballester2004np}, 
a coalition structure consisting of disjoint coalitions that cover all players, where each player estimates a valuation for other players in their group. Team formation has also been considered in project management, where evaluation of team members is conducted by measuring attributes such as leadership skills, technical talent, problem solving capabilities, and cultural relevance~\cite{%wi2009team, 
warhuus2021teaming}. Sports leagues assess the physical and functional well-being of players before forming teams~\cite{dadelo2014multi}. Such a task can be complex as %in competitive sports where 
teams often require players who are efficient in multiple roles. %Ahmed \textit{et al.}
\cite{ahmed2013multi} used the NSGA-II algorithm and employed an evolutionary multi-objective approach to obtain a high-performing team for cricket tournaments. Some factors that the system takes into consideration are batting average and bowling performance of individual players, wicket-keeper's past performances, and other rule-based constraints. Depending on the team selection strategy, each of the above factors are ranked (or nulled) according to importance and the final solution set is obtained according to it. If all the factors are deemed equally important, a domination approach is applied instead, i.e., sorting teams according to non-dominating factors (e.g., based on total cost of team formation) and picking the best front solution (i.e., lowest cost team).

Closer to our setting, \cite{machado2019fair} proposed brute force and heuristic approaches to create team recommendations in multidisciplinary projects, where most suitable candidates are incrementally selected until project requirements are fulfilled. However, unlike our work, they do not use any metrics to quantitatively measure the effectiveness of recommended teams, and additionally assume that an individual member may completely fulfill one skill requirement. The complexity of making multi-criteria decisions and building successful teams also increases with the number of available candidates. Evaluating the goodness of every team permutation quickly becomes computationally prohibitive and NP-hard~\cite{roy2014exploiting}. The semantics of group recommendation algorithms were also defined by a means of a Consensus Function, where \textit{disagreement} amongst group members was introduced as a factor to influence the generated recommendations~\cite{ameryahia2009group}. Alternatively, methods have also been proposed to seek a certain level of \textit{agreement} amongst group members, where individuals preferences are iteratively brought closer, until a desired threshold is achieved~\cite{castro2015consensus}. Search-based optimization techniques were also explored using a queuing simulation model to maximize resource usage in software project management~\cite{di2011use}. Another application of team formation is in entrepreneurial domains~\cite{lazar2020entrepreneurial}, where entrepreneurs search for trustworthy partners \textit{and} investors when building new ventures. 

Existing literature systematically answers many questions regarding what, why, and how teams are formed~\cite{costa2020team, juarez2021comprehensive}.
%, takai2017towards, zainal2020team}. 
However, recruiting a group of experts to work towards a common goal does not guarantee that they will \textit{always} operate as a team. %~\cite{juarez2021comprehensive}. 

\textbf{Drawbacks of Single-Item Recommendation.} A common objective of many recommender algorithms~\cite{su2009survey, cremonesi2010performance} is to learn each individual's preferences through his interactions with a system, estimate his satisfaction for items he has not interacted with before, and return the top-\textit{k} items with highest estimated ratings. However, many single-item recommender systems suffer from known issues such as popularity bias, %~\cite{abdollahpouri2019managing}, 
where suggestions show an uneven distribution in recommending similar items, and the cold start problem, %~\cite{panda2022approaches}, 
a state where a new user wishes to interact with a system but there is not enough information about his preferences to make an accurate decision.

\textbf{Motivation.} We therefore expand on our prior work, ULTRA (\textbf{U}niversity-\textbf{L}ead \textbf{T}eam Builder from \textbf{R}FPs and \textbf{A}nalysis)~\cite{srivastava2022ultra}, and consider a group recommendation setting that promotes research collaboration using novel AI methods to recommend optimal teaming suggestions. Our work encourages multi-functional and interdisciplinary teams to form, and brings together members from various disciplines of work responsibility. 